% R11 en ../write_thesis_roadmap.md. Fuentes: packages/evaluation/planilla/rasgos.csv (solo
% lectura; definición operativa resumida, nombres como en las tablas 3.1 y 5.2),
% packages/evaluation/planilla/prompts.md (variantes Esencial, literales) y las tablas que genera
% asg-results en redaccion/tablas/ (2026-11-03).
\appendix
% babel fija \appendixname («Apéndice») al empezar el documento; aquí se cambia por «Anexo».
\renewcommand{\appendixname}{Anexo}

\chapter{Lista núcleo de rasgos}\label{app:nucleo}

La Tabla~\ref{tab:anexo-nucleo} da los 24 rasgos núcleo de la Sección~\ref{sec:lista-nucleo} con su
identificador en el catálogo, que es el que usan las tablas generadas por código del
Capítulo~\ref{chapter:experiments}. La letra indica la capa del catálogo
(Sección~\ref{sec:catalogo}): T, medidas del texto calculadas por código; R, la rúbrica hecha
contable, y X, los recuentos de los ejemplos del profesor Piad. Los recuentos que crecen con la longitud se
expresan por mil palabras («/1000 pal.»). La última columna resume qué se cuenta; la definición
operativa completa está en la planilla del catálogo (Sección~\ref{sec:herramientas}).

\begin{table}[p]
  \centering
  \caption{Rasgos núcleo con su identificador, criterio y unidad}
  \label{tab:anexo-nucleo}
  \footnotesize
  \setlength{\tabcolsep}{4pt}
  \begin{tabular}{@{}l>{\raggedright\arraybackslash}p{3.5cm}>{\raggedright\arraybackslash}p{1.8cm}%
      >{\raggedright\arraybackslash}p{1.7cm}>{\raggedright\arraybackslash}p{6.3cm}@{}}
    \toprule
    Id & Nombre & Criterio & Unidad & Qué se cuenta \\
    \midrule
    T01 & Palabras & --- & recuento & palabras de la prosa (por código) \\
    T06 & Proporción de diálogo & personajes & proporción & párrafos con diálogo (por código) \\
    T08 & Repetición de trigramas & creatividad & proporción & trigramas repetidos (por código) \\
    R01 & Promesas de apertura & trama & recuento & promesas que deja el primer capítulo \\
    R02 & Señales de progreso & trama & /1000 pal. & avances del tramo medio \\
    R03 & Preparación del pago final & trama & recuento & preparaciones que cosecha el final (dos citas) \\
    R04 & Ciclos «sí, pero» & trama & recuento & intentos que logran algo con un coste nuevo \\
    R05 & Ciclos «no, y además» & trama & recuento & intentos fallidos que empeoran la situación \\
    R09 & Proactividad del protagonista & personajes & /1000 pal. & acciones que inicia el protagonista \\
    R16 & Soluciones preparadas & trama & recuento & conflictos resueltos con algo ya visto (dos citas) \\
    R18 & Límites del mundo & creatividad & recuento & una regla del mundo impide algo \\
    R23 & Escenas con cambio & trama & proporción & escenas en que algo cambia \\
    R24 & Decisiones con coste & personajes & recuento & decisiones que cuestan algo real \\
    X02 & Cuota del personaje más presente & personajes & proporción & escenas del personaje más presente \\
    X03 & Reparto de la participación & personajes & 0--1 & entropía de la presencia (por código) \\
    X06 & Escenas con conflicto & trama & proporción & escenas con objetivos enfrentados \\
    X10 & Mentiras & personajes & recuento & engaños deliberados a otro personaje \\
    X13 & Secretos revelados & trama & recuento & secretos que salen a la luz \\
    X14 & Complicaciones & trama & recuento & obstáculos nuevos a un objetivo \\
    X15 & Giros & creatividad & recuento & lo que se daba por cierto cambia \\
    X16 & Sucesos azarosos o externos & trama & recuento & sucesos que no provoca nadie \\
    X17 & Usos de objetos & trama & /1000 pal. & usar, dar, tomar, esconder o mostrar \\
    X20 & Contradicciones & trama & recuento & contradicciones internas (dos citas) \\
    X29 & Reparto del diálogo & personajes & 0--1 & entropía de las réplicas (por código) \\
    \bottomrule
  \end{tabular}
\end{table}

\chapter{Peticiones de las obras de los experimentos}\label{app:obras}

Las tres obras de la matriz y de las líneas base (Sección~\ref{sec:corpus}) recibieron
literalmente estas peticiones, la variante Esencial del catálogo. La línea base de un solo
\textit{prompt} recibe el mismo texto.

\section*{Obra 01: fantasía canónica}

\begin{quote}\small
Escribe en español un relato de fantasía épica con perfil narrativo Esencial. Sir Aldren, un caballero
veterano atormentado por el fracaso de una misión anterior, debe entrar en una fortaleza levantada
sobre un volcán para rescatar a la princesa Elara de un dragón ancestral. Elara no debe ser una víctima
pasiva: debe investigar su cautiverio, tomar decisiones arriesgadas y contribuir de forma decisiva a su
propia liberación. El dragón debe tener una motivación comprensible relacionada con una antigua promesa
rota por el reino, y no ser simplemente un monstruo malvado. Desarrolla una cadena causal clara desde la
llegada del caballero hasta el enfrentamiento final; prepara con antelación cualquier objeto,
conocimiento o habilidad que resulte decisivo. Mantén la continuidad de lugares, heridas, información y
relaciones. Usa un tono aventurero y emotivo, incluye un dilema moral que obligue a Aldren a elegir
entre obedecer al rey y hacer lo correcto, y termina con un desenlace cerrado y esperanzador. Evita el
deus ex machina, las profecías que resuelven el conflicto por sí solas y las explicaciones sobre el
proceso de escritura.
\end{quote}

\section*{Obra 03: «La última luz del faro»}

\begin{quote}\small
Escribe en español un relato de misterio con perfil narrativo Esencial. Durante una tormenta que deja
incomunicada una pequeña isla, la archivista Mara Vela llega al faro para catalogar sus registros y
descubre que el farero ha desaparecido de una habitación cerrada por dentro. En el edificio permanecen
su hija, un meteorólogo, la médica de la isla y un contrabandista retirado; todos ocultan algo, pero no
todos mienten sobre la desaparición. Construye un misterio de juego limpio: presenta antes de la
revelación todas las pistas necesarias para deducir qué ocurrió, incluida una anotación alterada, una
pieza del mecanismo del faro y una contradicción horaria. Distingue claramente lo que sabe cada
personaje y mantén una cronología consistente durante la tormenta. Incluye al menos una pista falsa que
tenga una explicación causal y no sea un engaño del narrador. No uses causas sobrenaturales, gemelos
secretos, amnesia ni confesiones que sustituyan la investigación. Usa una atmósfera tensa y aislada,
permite que Mara resuelva el caso relacionando evidencias observables y termina revelando tanto el
método como el motivo y las consecuencias humanas.
\end{quote}

\section*{Obra 04: «El cine de los domingos»}

\begin{quote}\small
Escribe en español un drama contemporáneo con perfil narrativo Esencial. Tras la muerte de su madre,
los hermanos Lucía y Tomás heredan un viejo cine de barrio que será demolido en siete días si no pagan
una deuda. Lucía quiere venderlo y regresar a la ciudad donde construyó su carrera; Tomás quiere
organizar una última función para demostrar que el lugar todavía importa. Ambos recuerdan de manera
diferente el abandono de su padre y creen que la madre favoreció al otro. Desarrolla el conflicto
mediante conversaciones, silencios, acciones prácticas y decisiones económicas concretas, sin
convertir a ninguno de los dos en villano. Un rollo de película incompleto, una libreta de cuentas y la
antigua cabina de proyección deben adquirir significado dramático y participar en la resolución. Haz
que cada cambio emocional tenga una causa visible y que la verdad sobre la familia complique el
conflicto en lugar de resolverlo de inmediato. Usa un tono íntimo y contenido, evita accidentes
oportunistas, enfermedades repentinas y herencias secretas, y termina con una decisión conjunta creíble
que implique una pérdida real y una forma limitada de reconciliación.
\end{quote}

\chapter{Participantes y alcance de las lecturas}\label{app:lectores}

La Tabla~\ref{tab:lectores} recoge los ocho participantes descritos por el autor en la actualización
del 7 de noviembre de 2026. Todos son graduados de Ciencia de la Computación. Las categorías de
lectura alta, media y baja son declaradas; no se les asignan cantidades de libros o horas que no
se registraron. L1 identifica al autor en este conjunto y no se corresponde necesariamente con el
L1 de las bases históricas del piloto.

\begin{table}[htb]
  \centering
  \caption{Perfiles anónimos del estudio (octubre--1 de noviembre de 2026)}
  \label{tab:lectores}
  \small
  % Revisión documental 2026-11-07; fuentes en revision.json.
\begin{tabular}{@{}llll@{}}
\toprule
Lector & Formación & Lectura & Conoce el sistema / autor \\
\midrule
L1 & Ciencia de la Computación & Alta & Sí / sí \\
L2 & Ciencia de la Computación & Alta & No / no \\
L3 & Ciencia de la Computación & Alta & No / no \\
L4 & Ciencia de la Computación & Media & No / no \\
L5 & Ciencia de la Computación & Media & No / no \\
L6 & Ciencia de la Computación & Media & No / no \\
L7 & Ciencia de la Computación & Baja & No / no \\
L8 & Ciencia de la Computación & Baja & No / no \\
\bottomrule
\end{tabular}

\end{table}

Cada lector recibió diez textos distintos en los mismos cinco pares y respondió tres preguntas por
par. Algunos realizaron además lecturas o validaciones, sin incorporar participantes nuevos. La
lista de títulos y ejecuciones, el detalle de los pares y las respuestas individuales no se han
identificado en los registros disponibles. Por ello, los códigos de las seis historias del piloto
anterior no se reutilizan para completar esta lista ni se atribuyen sus votos a estos lectores.

\chapter{Tablas complementarias de los experimentos}\label{app:tablas}

La Tabla~\ref{tab:anexo-matriz} conserva el estado histórico de la matriz según el manifiesto
actualizado el 31 de octubre de 2026 en UTC. Sus 24 unidades completadas forman el conjunto descrito
en la Sección~\ref{sec:resultados-horizontal}; el manifiesto también conserva una unidad fallida y
veinte pendientes del diseño inicial. Las nueve líneas base y la base simulada de la obra 03 ya
existían cuando se generaron las otras catorce unidades. La lista de ejecuciones y la comprobación de
extracciones se conservan en \codigo{redaccion/tablas/horizontal_revision.csv}.

\begin{table}[htb]
  \centering
  \caption{Estado histórico de la matriz horizontal (manifiesto, 2026-10-31)}
  \label{tab:anexo-matriz}
  \small
  \begin{tabular}{@{}lrrrr@{}}
\toprule
Celda & Unidades & Completadas & Fallidas & Sin generar \\
\midrule
M01 Base simulada & 6 & 5 & 1 & 0 \\
M02 Líneas base & 9 & 9 & 0 & 0 \\
M03 Narrativa desde el plan & 6 & 2 & 0 & 4 \\
M04 Tono sombrío & 6 & 2 & 0 & 4 \\
M05 Tono humorístico & 6 & 2 & 0 & 4 \\
M06 Visión limitada & 6 & 2 & 0 & 4 \\
M07 Primera persona & 6 & 2 & 0 & 4 \\
Total & 45 & 24 & 1 & 20 \\
\bottomrule
\end{tabular}

\end{table}

La Tabla~\ref{tab:extractor-detalle} amplía la validación registrada del extractor. Distingue los
veredictos conservados de los pares numéricamente comparables y los cambios de estado. Los rasgos
calculados por código no tienen veredictos sobre citas propios. También se muestran cuatro rasgos
que quedaron fuera del núcleo vigente, para no ocultar su denominador al resumir la validación.

\begin{table}[p]
  \centering
  \caption{Detalle de las hojas de revisión y del test-retest (2026-10-29)}
  \label{tab:extractor-detalle}
  \small
  % Revisión documental 2026-11-07; fuentes en revision.json.
\begin{tabular}{@{}llrrrrrr@{}}
\toprule
Id & Grupo & C & I & SR & Pares & $\Delta_r$ & Cambios \\
\midrule
T01 & código & --- & --- & --- & 6/6 & 0,00 & 0/6 \\
T06 & código & --- & --- & --- & 6/6 & 0,00 & 0/6 \\
T08 & código & --- & --- & --- & 6/6 & 0,00 & 0/6 \\
R01 & núcleo & 5 & 0 & 0 & 5/6 & 0,00 & 0/6 \\
R02 & núcleo & 15 & 0 & 3 & 4/6 & 0,00 & 0/6 \\
R03 & núcleo & 1 & 0 & 0 & 1/6 & 0,00 & 3/6 \\
R04 & núcleo & 5 & 0 & 2 & 2/6 & 0,00 & 1/6 \\
R05 & núcleo & 5 & 0 & 0 & 2/6 & 0,25 & 1/6 \\
R09 & núcleo & 10 & 0 & 6 & 4/6 & 0,19 & 0/6 \\
R16 & núcleo & 5 & 0 & 1 & 1/6 & 0,00 & 2/6 \\
R18 & núcleo & 6 & 0 & 7 & 4/6 & 0,12 & 1/6 \\
R23 & núcleo & 7 & 0 & 1 & 2/6 & 0,00 & 2/6 \\
R24 & núcleo & 8 & 0 & 0 & 2/6 & 0,00 & 2/6 \\
X02 & código & --- & --- & --- & 5/6 & 0,04 & 0/6 \\
X03 & código & --- & --- & --- & 5/6 & 0,01 & 0/6 \\
X06 & núcleo & 28 & 0 & 6 & 3/6 & 0,00 & 2/6 \\
X10 & núcleo & 0 & 0 & 1 & 1/6 & 0,00 & 2/6 \\
X13 & núcleo & 3 & 0 & 1 & 3/6 & 0,00 & 2/6 \\
X14 & núcleo & 21 & 0 & 1 & 4/6 & 0,00 & 1/6 \\
X15 & núcleo & 9 & 0 & 0 & 3/6 & 0,17 & 1/6 \\
X16 & núcleo & 6 & 0 & 0 & 4/6 & 0,18 & 0/6 \\
X17 & núcleo & 21 & 0 & 3 & 4/6 & 0,22 & 1/6 \\
X20 & núcleo & 1 & 0 & 1 & 1/6 & 0,00 & 2/6 \\
X29 & código & --- & --- & --- & 3/6 & 0,01 & 2/6 \\
R15 & sec. & 4 & 0 & 0 & 0/6 & --- & 3/6 \\
R17 & sec. & 0 & 0 & 0 & 1/6 & 0,00 & 1/6 \\
X09 & sec. & 0 & 0 & 0 & 0/6 & --- & 1/6 \\
X24 & sec. & 0 & 0 & 0 & 0/6 & --- & 2/6 \\
\bottomrule
\end{tabular}

  \par\medskip
  \begin{minipage}{0.96\textwidth}
    \footnotesize
    C: correctas; I: incorrectas; SR: sin veredicto resuelto (etiqueta histórica «dudoso»), excluidas
    de la precisión principal. Pares: historias medidas en ambas pasadas, de seis repetidas.
    $\Delta_r$: diferencia relativa media sobre esos pares. Cambios: transiciones de estado entre
    pasadas, de seis historias. El guion indica ausencia de medida o de aplicación; no es cero.
    Los campos de omisiones están vacíos y no permiten estimar cobertura. Los identificadores se
    definen en la Tabla~\ref{tab:anexo-nucleo} o en el catálogo para los rasgos secundarios.
  \end{minipage}
\end{table}

Las tablas y la figura del juez del piloto anterior se conservan en los archivos del proyecto,
con su manifiesto de fuentes del 3 de noviembre. No se incluyen como resultados del estudio
actualizado ni del corpus de 122 historias. La revisión documental identifica sus fuentes y los
nuevos resúmenes en \codigo{redaccion/tablas/revision.json}.
